% 31
\begin{frame}{Почему такая скорость может быть медленной}
\small
Если хотим улучшить точность в десять раз,
\[
\varepsilon
\longrightarrow
\frac{\varepsilon}{10},
\]
то верхняя оценка числа итераций увеличивается примерно в десять раз.

\begin{exampleblock}{Пример}
Переход от ошибки $10^{-2}$ к $10^{-3}$ в худшей гарантии требует на порядок большего $k$.
\end{exampleblock}
\end{frame}

% 32
\begin{frame}{Важно: это верхняя оценка, а не точный прогноз}
\small
Теорема утверждает
\[
f(x_k)-f^\star
\le
\frac{C}{k}.
\]
Она \textbf{не} утверждает, что
\[
f(x_k)-f^\star=\frac{C}{k}.
\]

\begin{alertblock}{Как читать оценки сходимости}
Они гарантируют: алгоритм работает \textbf{не хуже определённого сценария} на всём рассматриваемом классе функций. На конкретной задаче он может быть гораздо быстрее.
\end{alertblock}
\end{frame}

% 33
\begin{frame}{Теперь добавим сильную выпуклость}
\small
Пусть $f$ одновременно
\begin{itemize}
\item $L$-гладкая;
\item $\mu$-сильно выпуклая.
\end{itemize}
Тогда
\[
0<\mu\le L.
\]
Сильная выпуклость даст связь между величиной градиента и ошибкой по функции — именно её не хватало раньше.
\end{frame}

% 34
\begin{frame}{Что именно даёт сильная выпуклость}
\small
Для любого $y$:
\[
f(y)
\ge
f(x)+\grad f(x)^\top(y-x)
+\frac\mu2\norm{y-x}^2.
\]
Правая часть — квадратичная функция по $y$.

\begin{block}{Идея}
Если она является глобальной нижней оценкой $f$, можно минимизировать её и получить нижнюю оценку на само $f^\star$.
\end{block}
\end{frame}

% 35
\begin{frame}{Минимизируем квадратичную нижнюю модель}
\small
Положим
\[
h=y-x.
\]
Нужно минимизировать по $h$
\[
q(h)=\grad f(x)^\top h+\frac\mu2\norm{h}^2.
\]
Условие стационарности:
\[
\grad_h q(h)=\grad f(x)+\mu h=0.
\]
Следовательно,
\[
\boxed{
h^\star=-\frac1\mu\grad f(x).
}
\]
\end{frame}

% 36
\begin{frame}{Минимальное значение модели}
\small
Подставляем $h^\star=-\frac1\mu\grad f(x)$:
\[
\begin{aligned}
q(h^\star)
&=-\frac1\mu\norm{\grad f(x)}^2
+\frac\mu2\frac1{\mu^2}\norm{\grad f(x)}^2\\
&=-\frac1{2\mu}\norm{\grad f(x)}^2.
\end{aligned}
\]
Поэтому для любого $y$ и, в частности, для минимума:
\[
f^\star
\ge
f(x)-\frac1{2\mu}\norm{\grad f(x)}^2.
\]
\end{frame}

% 37
\begin{frame}{Сильная выпуклость не позволяет градиенту быть слишком маленьким}
\small
Переносим члены:
\[
 f(x)-f^\star
 \le
 \frac1{2\mu}\norm{\grad f(x)}^2.
\]
Или эквивалентно
\[
\boxed{
\norm{\grad f(x)}^2
\ge
2\mu\bigl(f(x)-f^\star\bigr).
}
\]
Это ключевой мост между леммой об убывании и ошибкой по функции.
\end{frame}

% 38
\begin{frame}{Геометрический смысл предыдущего неравенства}
\centering
\includegraphics[width=0.82\linewidth]{gradient_gap_relation.pdf}

\vspace{0.15em}
\small
В сильно выпуклой задаче заметное отставание по функции обязано сопровождаться достаточно большим градиентом.
\end{frame}

% 39
\begin{frame}{Снова используем лемму об убывании}
\small
При шаге $1/L$:
\[
\gap{x_{k+1}}
\le
\gap{x_k}
-
\frac1{2L}\norm{\grad f(x_k)}^2.
\]
Из сильной выпуклости:
\[
\norm{\grad f(x_k)}^2
\ge
2\mu\bigl(f(x_k)-f^\star\bigr).
\]
Подставляем нижнюю оценку на градиент в отрицательный член.
\end{frame}

% 40
\begin{frame}{Получаем сокращение ошибки за один шаг}
\small
\[
\begin{aligned}
\gap{x_{k+1}}
&\le
\gap{x_k}
-
\frac1{2L}\cdot2\mu\bigl(f(x_k)-f^\star\bigr)\\
&=
\left(1-\frac\mu L\right)\bigl(f(x_k)-f^\star\bigr).
\end{aligned}
\]
То есть
\[
\boxed{
\gap{x_{k+1}}
\le
\left(1-\frac\mu L\right)\bigl(f(x_k)-f^\star\bigr).
}
\]
Каждая итерация сокращает верхнюю оценку на один и тот же множитель.
\end{frame}

% 41
\begin{frame}{Повторяем неравенство $k$ раз}
\small
Последовательно:
\[
\gap{x_1}
\le
q\bigl(f(x_0)-f^\star\bigr),
\]
\[
\gap{x_2}
\le
q\bigl(f(x_1)-f^\star\bigr)
\le
q^2\bigl(f(x_0)-f^\star\bigr),
\]
где
\[
q=1-\frac\mu L.
\]
Поэтому
\[
\boxed{
\gap{x_k}
\le
\left(1-\frac\mu L\right)^k\bigl(f(x_0)-f^\star\bigr).
}
\]
\end{frame}

% 42
\begin{frame}{Теорема о линейной сходимости}
\small
\begin{block}{Теорема}
Пусть $f$ $\mu$-сильно выпукла и $L$-гладка. Для GD с шагом $1/L$:
\[
\boxed{
 f(x_k)-f^\star
 \le
 \left(1-\frac\mu L\right)^k
 \bigl(f(x_0)-f^\star\bigr).
}
\]
\end{block}

\begin{alertblock}{Терминология}
«Линейная сходимость» не означает движение по прямой. Это стандартное название геометрического закона $e_{k+1}\le q e_k$, где $q<1$.
\end{alertblock}
\end{frame}

% 43
\begin{frame}{$1/k$ против геометрической скорости}
\centering
\includegraphics[width=0.88\linewidth]{rate_preview.pdf}

\vspace{0.15em}
\small
Геометрическое уменьшение особенно хорошо видно на логарифмической шкале: оно становится почти прямой линией.
\end{frame}

% 44
\begin{frame}{Число обусловленности снова появляется}
\small
Определим
\[
\boxed{\kappa=\frac{L}{\mu}.}
\]
Тогда множитель сокращения
\[
q
=1-\frac\mu L
=1-\frac1\kappa.
\]
Следовательно,
\[
\boxed{
\gap{x_k}
\le
\left(1-\frac1\kappa\right)^k\bigl(f(x_0)-f^\star\bigr).
}
\]
\end{frame}

% 45
\begin{frame}{Хорошо обусловленная задача}
\small
Пусть
\[
\kappa=2.
\]
Тогда
\[
q=1-\frac12=0.5.
\]
То есть гарантированная верхняя оценка ошибки делится примерно пополам на каждой итерации:
\[
1,\;0.5,\;0.25,\;0.125,\ldots
\]
\end{frame}
